\section{Valuation Representation}
\label{sec:representation}

Every valuation on a finite frame is a probability mass function on
points. On infinite frames, this can fail, and Scott continuity
marks the boundary between where it still holds and where
it does not.

\subsection{Finite Frames Representation}

\begin{theorem}[Birkhoff duality]
Every finite distributive lattice is isomorphic to the lattice of
lower sets of a finite poset: its sup-irreducible elements form a
poset $P$ whose lower sets reconstruct the lattice, and every finite
poset arises this way from its own lattice of lower sets.
\end{theorem}

\noindent
Let $P$ be the finite poset of
sup-irreducible elements of a finite frame.
Then, by Birkhoff duality
(\texttt{OrderIso.lowerSetSupIrred}),
every finite frame is the lattice of \emph{lower sets} of $P$: a
\emph{lower set} $U \subseteq P$ satisfies $p \in U,\, q \le p
\Rightarrow q \in U$, with $\mathrm{Iic}\,p = \{q : q \le p\}$ and
$\mathrm{Iio}\,p = \{q : q < p\}$ its principal instances at $p$.
We prove the finite representation theorem below on
\texttt{LowerSet P}
(\texttt{Valuation.\allowbreak exists\_pmf\_presentation}).  
Transporting
it across the isomorphism above then gives the theorem for every
finite frame, not just lower sets.  Let us define the \emph{mass} of a
point as the discrete derivative of $v$:
\begin{lstlisting}
def Valuation.mass (v : Valuation (LowerSet P)) (p : P) :
    ℝ≥0∞ := v (LowerSet.Iic p) - v (LowerSet.Iio p)
\end{lstlisting}

From the mass at every point, we can fully recover $v$.
The masses also sum to $1$, forming a
probability distribution on $P$:

\begin{theorem}[finite representation,
  \texttt{eq\_sum\_mass}, \texttt{sum\_mass}]
\label{thm:finite-rep}
For a finite $P$, any valuation $v$ on $\mathrm{Low}(P)$, and any
$U \in \mathrm{Low}(P)$,
$v(U) = \sum_{p \in U} \mathrm{mass}(p)$, and
$\sum_{p} \mathrm{mass}(p) = 1$.
\end{theorem}

The proof works by induction on $|U|$: repeatedly remove a maximal
point $p$ from $U$, and use modularity applied to
$(U \setminus {\uparrow}p,\; {\downarrow}p)$ to show each removal
changes $v$ by exactly $\mathrm{mass}(p)$.
\texttt{Valuation.toPMF} gives
$\mathrm{mass}$ the type \texttt{PMF} on $P$, i.e. a
classical probability mass function.  And since lower sets are
exactly the opens of the Alexandrov topology, where interior is the
identity, the theorem also reads $v = \mu \circ \mathrm{int}$, i.e.
every valuation on a
finite frame is the classical measure $\mu$ applied after
taking the interior.
In the finite case, this is the converse of
the measure-theoretic bridge of Section~\ref{sec:bridges}.

The theorem above represented $v$ as a sum of point masses.
The dual view
represents $v$ as a mixture of sharp
\emph{point-valuations}.  For each $p \in P$, let
$\delta_p$ be the valuation with $\delta_p(U) = 1$ if $p \in U$, and
$\delta_p(U) = 0$ otherwise:

\begin{theorem}[mixture characterization,
  \texttt{eq\_mix\_deltaPoint}]
\label{thm:mixture}
On a finite frame, every valuation $v$ satisfies
$v = \sum_p \mathrm{mass}(p)\cdot\delta_p$.
\end{theorem}

\noindent
Together with the fact that any mixture of valuations is again a
valuation (\texttt{Valuation.mix}), this theorem characterizes the
valuations on a finite frame exactly: they are the finite convex
mixtures of point-valuations, closed under mixing because modularity
is linear.  Section~\ref{sec:cox} shows modularity cannot instead be
\emph{derived} from more primitive axioms.  It must be posited
outright, and doing so is exactly what makes every valuation a
mixture of sharp, all-or-nothing certainties.

\subsection{Infinite Boundary Representation via Scott Continuity}

Both halves of the finite theorem fail for infinite frames:

\begin{theorem}[obstruction,
  \texttt{exists\_valuation\_not\_\allowbreak point\_representable}]
\label{thm:obstruction}
On $\mathrm{Low}(\mathbb{N})$, the indicator $v$ of $\top$ is a
valuation with $\mathrm{mass}(p) = 0$ for every point $p$, and
$v(\top) = 1$.
\end{theorem}

\noindent
No point carries any of the mass.  The failure is one of
approximation: $v(\top)$ is not the supremum of $v(\mathrm{Iic}\,n)$
over the finite stages. In order to recover the representation,
we introduce Scott continuity:

\begin{definition}[directed set, Scott continuity]
A subset $D \subseteq \Omega$ is \emph{directed} if it is nonempty
and every pair of its elements has an upper bound in $D$.  A
valuation $v$ is \emph{Scott-continuous} if
$v(\bigsqcup D) = \sup_{a \in D} v(a)$ for every directed $D$.
\end{definition}

We must require Scott continuity 
explicitly, as the representation result may fail otherwise. 
For example, Scott continuity (\texttt{topIndicator\_not\_scott}) 
failure can be demonstrated for a whole class of valuations
(Theorem \ref{thm:recovery}):

\begin{theorem}[recovery, \texttt{eq\_tsum\_mass\_of\_scott},
  generalized by \texttt{isPurelyAtomic\_of\_scott}]
\label{thm:recovery}
Let $v$ be a valuation on $\mathrm{Low}(\mathbb{N})$. If
$v(\top) = \bigsqcup_n v(\mathrm{Iic}\,n)$
then $v(\top) = \sum'_n \mathrm{mass}(n)$.
\end{theorem}

In general, on
any poset all of whose principal downsets are finite, Scott
continuity implies zero diffuse part without additional preconditions.
However, subtraction in $\ENNReal$ is truncated at $0$, so if
$\sum'_p \mathrm{mass}(p)$ were greater than $v(\top)$, the expression
$v(\top) - \sum'_p \mathrm{mass}(p)$ would still evaluate to $0$.
Theorem~\ref{thm:decomposition} rules out
$\sum'_p \mathrm{mass}(p) > v(\top)$ unconditionally, so this
subtraction is never actually truncated and its value is 
the actual result of the subtraction:

\begin{theorem}[general decomposition, \texttt{tsum\_mass\_le}]
\label{thm:decomposition}
For every valuation $v$ on $\mathrm{Low}(P)$ and an arbitrary poset $P$,
$\sum'_p \mathrm{mass}(p) \le v(\top)$.
\end{theorem}

\noindent
Rearranging Theorem~\ref{thm:decomposition},
licensed by the truncation point above, gives
$v(\top) - \sum'_p \mathrm{mass}(p) \ge 0$.  Every valuation thus
splits unconditionally into an \emph{atomic} part carried by points
($\sum'_p \mathrm{mass}(p)$, which is always a sub-probability), and a
\emph{diffuse} (nonnegative) remainder
$v(\top) - \sum'_p \mathrm{mass}(p)$.
This is a constructive Lebesgue-style decomposition.  The finite and
Scott cases are the purely atomic extremes, and the $\top$-indicator
is the purely diffuse one.  The proof of
Theorem~\ref{thm:decomposition} re-runs the maximal-point induction on
each finite subset of $P$, bounding its partial sum by $v(\top)$,
then takes the supremum over all finite subsets to bound
$\sum'_p \mathrm{mass}(p)$ itself.  
Its proof does not directly invoke an excluded-middle case split.
However, it relies on the (classical) decidability of the ordering
of the reals (a default in \textsf{mathlib}).
Representing the diffuse part on non-spatial frames remains the open
frontier.  The obstruction
there is again that the frame has \emph{too few points}, the same
one as behind the undecidability of Section~\ref{sec:bridges}.
